\documentclass[aspectratio=169]{beamer}


\usetheme[
  style=minimal,
  palette=horizon,
  frametitle=subtle,
  progressbar,
  sansface={libertinus},
%   mathfont={libertinus},
]{Celestia}

\title{Macroeconomics in the Sequence Space}
\author{Chen Gao \and Weigao Wang}
\institute{School of Economics, Peking University}
\date{\today}


\begin{document}

\maketitle

\begin{frame}{Outline}
    \tableofcontents
\end{frame}

\section{Sequence Space vs.\ State Space}

\begin{frame}{Sequence Space vs.\ State Space}{Motivation}

    \begin{itemize}
        \item In heterogeneous-agent (HA) models, the economy’s state is the \alert{entire cross-sectional distribution} — an infinite-dimensional object.
        \item Sequence-space methods \emph{sidestep} this challenge: they characterize \alert{local dynamics around the steady state} without ever carrying the distribution as a state variable.
        \item The key insight: aggregate equilibrium responses to shocks can be encoded in \alert{Jacobian matrices} of aggregate outcomes with respect to aggregate inputs.
    \end{itemize}

\end{frame}


\begin{frame}{Sequence Space vs.\ State Space}{State Space}

    \begin{definition}
        A \alert{state-space representation} describes the economy recursively.
        At date $t$, a state variable $x_t$ summarizes all payoff-relevant history, and the model evolves as
        \[
            x_{t+1}=f(x_t,\varepsilon_{t+1}),
            \qquad
            y_t=g(x_t).
        \]
        \emph{Given $x_t$, the future is independent of the past.}
    \end{definition}
\end{frame}

\begin{frame}{Sequence Space vs.\ State Space}{Sequence Space}

    \begin{definition}
        A \alert{sequence-space representation} characterizes equilibrium as a mapping between entire time paths. We seek sequences $\{X_t\}_{t=0}^{\infty}$ such that equilibrium conditions hold at every date:
        \[
            F\!\left(\{X_t\}_{t=0}^{\infty},\{Z_t\}_{t=0}^{\infty}\right)=\mathbf{0},
        \]
        where $\{X_t\}$ collects endogenous variables and $\{Z_t\}$ collects exogenous shocks or policy paths.
    \end{definition}

\end{frame}


\begin{frame}{Sequence Space vs.\ State Space}{Comparison}

    \begin{block}{Intuition}
        \begin{itemize}
            \item State space asks: \emph{what is the current state, and how does it transition?}
            \item Sequence space asks: \emph{given a shock path, what equilibrium outcome path clears all markets?}
        \end{itemize}
    \end{block}

    \begin{block}{Takeaway}
        When the distribution is the intractable state variable, sequence space reframes a ``huge state'' problem as a \alert{structured Jacobian} problem — amenable to fast, scalable computation.
    \end{block}
\end{frame}




\section{SSJ in Krusell-Smith Model}
\begin{frame}{Example: The Krusell-Smith Model}{Setup}
    \begin{itemize}
        \item The economy is populated by a mass 1 of heterogeneous households,
        \item Each household maximizes its expected lifetime utility:
              \[
                  \mathbb{E}_0\left[\sum_{t=0}^{\infty} \beta^t u(c_t)\right],
              \]
        \item Standard CRRA utility,
        \item There exist $n_e$ idiosyncratic states w/ transition matrix $P(e,e')$ and stationary distribution $\pi$,
        \item Agents supply an exogenous number of hours $n$, and earn wage income $w_t e n$,
        \item Agents can only trade in capital $k$, which pays a rental rate $r_t$ net of depreciation.
    \end{itemize}


\end{frame}

\begin{frame}{Example: The Krusell-Smith Model}{Household Problem}
    The value function of an agent entering period w/ state $(e,k_-)$ is given by:

    \begin{align*}
        V_t(e, k_-)               & = \max_{c,\,k}  u(c) + \beta \sum_{e'} V_{t+1}(e', k) \, P(e,e') \\
        \text{s.t.} \qquad  c + k & = (1 + r_t)k_- + w_t e n                                         \\
        k                         & \geq 0
    \end{align*}

    \begin{itemize}
        \item Denote by $c^*_t(e, k_-)$ and $k^*_t(e, k_-)$ the policy functions that solve the Bellman equation.
        \item Also denote by $D_t(e, K_-)$ the measure of households in state $e$ that own capital in a set $K_-$ at the start of date $t$.
    \end{itemize}
\end{frame}


\begin{frame}{Example: The Krusell-Smith Model}{Household Problem}
    The distribution $D_t$ has law of motion
    \[
        D_{t+1}(e', K) = D_t(e, k^{*-1}_t(e, K)) P(e,e').
    \]
    where $k^{*-1}_t(e, K)$ denotes the inverse of $k^*_t(e, K)$.

    \begin{itemize}
        \item We assume that prior to $t = 0$, the economy is in a steady state.
        \item We suppose that agents start in this stationary distribution at date $0$, so that $D_0 = D_{ss}$.
    \end{itemize}

\end{frame}


\begin{frame}{Example: The Krusell-Smith Model}{Aggregate Capital}
    \begin{itemize}
        \item For any $t$, the policy $k^*_t(e, k_-)$ is a function of the future path
              $\{r_s, w_s\}_{s=t}^{\infty}$.
        \item \alert{Given $D_0 = D_{ss}$, the distribution $D_t(e, K)$ at any $t$ is a function of the entire path $\{r_s, w_s\}_{s=0}^{\infty}$.}
        \item The aggregate household capital holdings are characterized by a capital function $\mathcal{K}_t(\{r_s, w_s\}_{s=0}^{\infty})$.

              \[
                  \mathcal{K}_t\left(\{r_s, w_s\}_{s=0}^{\infty}\right) =
                  \sum_{e} \int_{k_-}
                  k^*_t(e, k_-) D_t(e, dk_-)
              \]
    \end{itemize}

    \begin{alertblock}{Key Insight}
        The ability to reduce interactions between heterogeneous agents to functions such as $\mathcal{K}_t$,
        which map aggregate sequences into aggregate sequences, is key to the sequence-space Jacobian method.
    \end{alertblock}
\end{frame}


\begin{frame}{Example: The Krusell-Smith Model}{Production}
    \begin{itemize}
        \item Firm adopts a Cobb-Douglas production function: \[Y_t = Z_t K_t^{\alpha} N_t^{1-\alpha}\]
        \item The firm’s first-order conditions:
              \begin{align*}
                  r_t & = \alpha Z_t K_{t-1}^{\alpha-1} N_t^{1-\alpha} -\delta \\
                  w_t & = (1 - \alpha) Z_t K_{t-1}^{\alpha} N_t^{-\alpha}
              \end{align*}

    \end{itemize}
\end{frame}
\section{SSJ development}

\begin{frame}{SSJ in Continuous Time}
    \textit{To be developed}
\end{frame}

\begin{frame}{Determinacy and Large-Scale Solutions in the Sequence Space}
    \textit{To be developed}
\end{frame}

\begin{frame}{2nd SSJ}
    \textit{To be developed}
\end{frame}

\subsection{Application}

\begin{frame}{fiscal and monetary policy}
    \textit{To be developed}
\end{frame}

\begin{frame}{behavioural macro}
    \textit{To be developed}
\end{frame}

\begin{frame}{life-cycle model}
    \textit{To be developed}
\end{frame}

\begin{frame}{dynamic spatial model}
    \textit{To be developed}
\end{frame}

\begin{frame}{production network}
    \textit{To be developed}
\end{frame}

\begin{frame}{het households + het firms}
    \textit{To be developed}
\end{frame}

\begin{frame}{het banks}
    \textit{To be developed}
\end{frame}

\begin{frame}{optimal policy}
    \textit{To be developed}
\end{frame}

\begin{frame}{macro labor}
    \textit{To be developed}
\end{frame}

\section{Global Solution}

\subsection{Global Nonlinear Solutions in Sequence Space and the Generalized Transition Function}

\begin{frame}{Global Nonlinear Solutions in Sequence Space and the Generalized Transition Function}
    \textit{To be developed}
\end{frame}

\subsection{DL in sequence space}

\begin{frame}{DL in sequence space}
    \textit{To be developed}
\end{frame}

\subsection{RL in sequence space}

\begin{frame}{RL in sequence space}
    \textit{To be developed}
\end{frame}

\section{Reference}

\begin{frame}{Reference}
    \textit{To be developed}
\end{frame}

\end{document}
